\section{Conclusion}

We benchmarked the LSODA, BDF, Radau, and CVODE solvers on the 1D p-Laplacian equation. The results suggest that utilizing a sparse Jacobian structure is generally necessary to avoid \(\mathcal{O}(N_x^3)\) scaling walls at high spatial resolutions. Across the tested parameter space with \(p > 2\), LSODA provided a strong balance of speed and precision, seemingly managing problem stiffness more effectively than BDF or Radau. CVODE appears to be a highly competitive alternative that may outperform LSODA in strongly degenerate regimes (e.g., \(p > 3\)), though its strictly implicit architecture leaves it potentially vulnerable to catastrophic convergence failures at extreme nonlinearities without careful configuration. Our extreme spatial scaling tests revealed that even at massive grid resolutions (\(N_x = 100,000\)), \pkg{SciPy}'s LSODA remains highly competitive with CVODE's optimized C-level backend.

However, when we tested the singular regime (\(p < 2\)), the solver dynamics shifted radically. As the regularization parameter \(\varepsilon \to 0\), the system generates an infinite initial diffusion transient. Under these conditions, the BDF implementations within LSODA and CVODE encounter catastrophic linear solver collapse due to extreme Jacobian ill-conditioning that rapidly exceeds double-precision floating-point limits. In contrast, the strictly L-stable architecture of the Radau solver successfully navigates this regime by heavily damping the infinite transient. Nonetheless, this stability comes at an extreme computational cost as \(N_x\) becomes large.

On the basis of our tests, no single solver is universally reliable for this PDE across all nonlinearity indices. For degenerate cases (\(p > 2\)), LSODA seems highly robust; its automatic method-switching architecture adapts to the problem's evolving stiffness. Conversely, for strongly singular cases (e.g., \(p < 1.01\)), fully implicit L-stable methods like Radau are required to gracefully navigate initial-condition singularities that otherwise cause catastrophic failures in standard stiff solvers. We therefore conclude that the numerical integration of the \(p\)-Laplacian equation resists the paradigm of a universal solver. The dichotomy between its degenerate and singular limits can cause algorithmic advantages in one regime to become liabilities in the other, and therefore, there are no definitive black-box solutions. 

The benchmarks also highlighted several notable solver behaviors. Loose tolerances appeared to induce severe performance anomalies and internal thrashing in both Radau and CVODE. Additionally, we observed that in the degenerate regime (\(p > 2\)), shrinking the regularization parameter (\(\varepsilon\)) seems to increase system stiffness only up to a threshold of roughly \(10^{-6}\). Beyond this point, further reductions did not induce significant additional computational overhead in our tests. Note that this plateau does not apply to the singular regime (\(p < 2\)), where further reductions exponentially amplify stiffness to the point of arithmetic failure.

We further analyzed the impact of spatial discretization methods and whether providing analytically derived Jacobians to time integrators improves stability and speed. In our tests, the finite-difference approach was twice as fast as using a variational (weak) formulation via finite elements. This is expected in 1D due to the inherent computational overhead of finite elements, though the finite-element method remains advantageous for its scalability to higher dimensions and ease of constructing analytic Jacobians. Furthermore, providing analytically derived Jacobians did not yield universal performance gains. While BDF and Radau from the \pkg{SciPy} suite became two to three times faster, LSODA's performance improved only marginally with increasing stiffness. Conversely, CVODE became significantly slower, and the cause is unclear. However, in the extreme singular limit (\(\varepsilon = 10^{-30}\)), providing the analytical Jacobian unexpectedly rescued CVODE from finite-difference conditioning errors that otherwise caused it to fail.

Finally, these performance profiles remain highly contingent on the host architecture: memory latency, cache topologies, and linear‑algebra backends can all shift the relative results. The entire benchmark suite is libre software and deliberately modular, so we warmly invite other researchers not only to reproduce the experiments on their own hardware, but to remix the components, explore untested parameter corners, and discover new phenomena that might have escaped our sweep.
